\section{Separation of two solutions with the same datum}\label{sec:separation}

We now combine the two compactness procedures. All parameters have been
fixed as in Proposition~\ref{prop:finite-dilated-horizon}, with the final
small parameters chosen to satisfy the preceding cold and cap estimates.
Fix one dormant limit \(F^{\mathrm d}\) on \([0,T_0]\), supplied by
Proposition~\ref{prop:admissible-limits}.

\begin{lemma}\label{sep:dormant-vanishing}
As \(T\downarrow0\),
\[
T^{-1}F^{\mathrm d}(Tt,Tx,v)\longrightarrow0
\]
in distributions on compact subsets of
\[
0<t<1,\qquad x\in\R^3,\quad v\in\R^3,\quad \bar v\ne0 .
\]
The evaluation in \(x\) uses the periodic lift centered at the origin.
\end{lemma}

\begin{proof}
The dormant majorant passes to the fixed-horizon weak limit. Its hot
part is at most \(t_{\rm old}^2H\), so after dilation and multiplication
by \(T^{-1}\) it is bounded by \(Tt^2H\), which vanishes in \(L^1\) on
the indicated compact sets.

For the cold majorant, a main enlarged annulus at physical transverse
position \(y_{\rm old}=Ty\), with \(y\) bounded, has \(s\le C T\).
The relevant periodic copy is the one centered at the origin for all
sufficiently small \(T\); other copies contribute only escaped tails.
For \(|\bar v|\ge\varepsilon>0\), the transverse Gaussian in the
cold estimate gives, after enlarging fixed constants,
\[
c_{\rm main}(Tt,Tx,v)
 \le C s^{-31}\exp\!\left(-\frac{c\varepsilon^2}{d_0^2s^{20}}\right).
\]
The right side, maximized over \(0<s\le CT\), tends to zero faster
than any power of \(T\). At points with no main color only tails
are present. The escaped sum is at most
\(C_A(Tt)^Ae^{-4|v|^2}\), and after multiplication by \(T^{-1}\)
it vanishes for \(A>1\). These uniform bounds on the compact set
prove the claim.
\end{proof}

\begin{proof}[Proof of Theorem~\ref{thm:main}]
The datum is \eqref{eq:datum}; Lemma~\ref{lem:datum} verifies all its
integrability requirements and bounded velocity support.
For every horizon \(0<T\le T_0\), Proposition~\ref{prop:finite-cap}
gives cap approximations. Proposition~\ref{prop:admissible-limits}
gives at least one admissible limit on \([0,T]\), with the same datum
as \(F^{\mathrm d}\).

Suppose, to obtain a contradiction, that there are horizons
\(T_k\downarrow0\) at each of which a cap limit can equal
\(F^{\mathrm d}\) on that interval. At each fixed \(T_k\), choose
approximations along a sequence converging to this particular common
limit. Use a countable exhaustion by nonnegative smooth compact tests
of the open region in Lemma~\ref{sep:dormant-vanishing}. A diagonal
choice of the approximation indices \(N_k\) makes \(s_{N_k}/T_k\to0\)
and gives
\begin{equation}\label{sep:diagonal}
f_k(t,x,v):=T_k^{-1}F_k(T_kt,T_kx,v)\longrightarrow0
\end{equation}
distributionally throughout that region.
Here \(F_k\) denotes the selected finite approximation.
For completeness, a fixed rescaled test at a fixed \(T_k\) is simply
a smooth physical test with finite, possibly large norm. Weak
convergence at that horizon permits its error to be made arbitrarily
small. One then enforces finitely many such errors and a sufficiently
small \(s_N/T_k\) simultaneously. No convergence rate uniform in
\(T_k\) is being assumed.

Choose a cap event of \(F_k=c_k+h_k\) and write
\(r_k=t_{*,k}+s_{N_k}\). Proposition~\ref{prop:blowup-limit} gives,
after subselection,
\[
\frac{s_{N_k}}{r_k}\to l,\qquad \frac{T_k}{r_k}\to U\in[1,\infty],
\]
and a dilated hot limit \(g\) on \(0<t<U\).
Proposition~\ref{prop:nonzero-limit} says that \(g\) is nonzero.
This statement supplies an observable lower bound. Since \(g\) is a
nonnegative density and the velocity axis has zero measure, there are
\(\psi\in C_c^\infty((0,U)\times\R^3_x\times\{\bar v\ne0\})\),
\(0\le\psi\le1\), and \(\delta>0\) such that
\[
 \int\psi g>2\delta,\qquad
 \int\psi(t,x,v)\,r_k^{-1}h_k(r_kt,r_kx,v)\dd t\dd x\dd v
 \ge\delta
\]
for all sufficiently large \(k\). The second inequality follows from
the local weak convergence in Proposition~\ref{prop:blowup-limit}.
The test and lower bound may depend on the extracted limit; no uniform
numerical value is needed.
Proposition~\ref{prop:finite-dilated-horizon} says that \(U<\infty\).
Consequently \(\rho_k=r_k/T_k\to\rho=1/U>0\).

Since \(F_k\ge h_k\),
\[
0\le r_k^{-1}h_k(r_kt,r_kx,v)
 \le\rho_k^{-1}f_k(\rho_kt,\rho_kx,v).
\]
For any compact set with \(0<t<U\), bounded \(x,v\), and
\(\bar v\ne0\), its image under the right-hand dilation lies in one
fixed compact subset of the region in \eqref{sep:diagonal}, for all
large \(k\). Positivity allows a single smooth nonnegative test,
equal to one on this larger compact set, to dominate the changing
tests. The right side therefore tends to zero in local \(L^1\).
It follows that \(g=0\) off the velocity axis. Since
\(0\le g\le(t+l)H\) is a density and the axis has Lebesgue measure
zero, \(g=0\) almost everywhere. This contradicts
Proposition~\ref{prop:nonzero-limit}.

Thus for all sufficiently small horizons no cap limit can equal
the fixed dormant limit. Choose one such horizon and one cap limit.
The two early solutions have the same datum and differ on a set of
positive space-time-velocity measure. Proposition~\ref{prop:admissible-limits}
continues both globally, preserving Definition~\ref{def:admissible}.
Their early difference persists, which proves the theorem.
\end{proof}

The separation occurs at finite velocities near the initial spatial
singularity. The conclusion requires neither an incoming flux from
infinite velocity nor an alteration of the conservation and entropy
conditions.
